\section{Experiments}

\subsection{RQ1: How well can execution-free methods tell whether the model's own verdict can be trusted? }
\label{sec:rq-verdict}

% Generated by scripts/paper/make_result_tables.py -- do not edit.
\begin{table}[t]
\centering\footnotesize
\setlength{\tabcolsep}{3.2pt}
\caption{AUROC for verdict correctness (mean over 3 partitions). \emph{LCB}, \emph{BCB}: trained and evaluated on that benchmark; $\to$: evaluated on the other one. Language shift: \emph{ID} is the mean over the three training languages, \emph{Shift} the mean over the six OOD train$\to$evaluation pairs (Table~\ref{tab:language-detail-verdict}), $\Delta$ their difference. %In-distribution cells are cross-fitted over problems; shifted cells use the layer chosen in distribution. 
Best per column and LLM in bold.}
\label{tab:verdict-summary}
\begin{tabular}{l rrrr rrr}
\toprule
 & \multicolumn{4}{c}{Benchmark shift} & \multicolumn{3}{c}{Language shift} \\
\cmidrule(lr){2-5}\cmidrule(lr){6-8}
Method & LCB & $\to$BCB & BCB & $\to$LCB & ID & Shift & $\Delta$ \\
\midrule
\multicolumn{8}{l}{\emph{Qwen3-Coder-30B}} \\
Confidence & .672 & \textbf{.552} & .555 & .669 & .644 & .643 & \textbf{$-$.001} \\
Probe & .699 & .480 & \textbf{.639} & .516 & \textbf{.761} & .668 & $-$.093 \\
\addlinespace[2pt]
Counterf.\ sim. & .692 & .453 & .560 & .680 & .759 & \textbf{.728} & $-$.031 \\
Value probe & \textbf{.740} & .485 & .579 & \textbf{.695} & .758 & .728 & $-$.031 \\
\midrule
\multicolumn{8}{l}{\emph{gpt-oss-20b}} \\
Confidence & .569 & .598 & .597 & .544 & .542 & .536 & \textbf{$-$.005} \\
Probe & \textbf{.829} & .478 & .507 & .476 & .774 & .668 & $-$.106 \\
\addlinespace[2pt]
Counterf.\ sim. & .779 & .572 & \textbf{.601} & .556 & \textbf{.808} & \textbf{.691} & $-$.116 \\
Value probe & .695 & \textbf{.599} & .600 & \textbf{.605} & .713 & .646 & $-$.067 \\
\bottomrule
\end{tabular}
\end{table}

% Generated by scripts/paper/make_result_tables.py -- do not edit.
\begin{table*}[t]
\centering\footnotesize
\setlength{\tabcolsep}{1.5pt}
\caption{AUROC for verdict correctness, in distribution (ID) and under language shift per training language (LCB problems; ID cross-fitted, shifted columns at the layer chosen in distribution). Mean $\pm$ SD over 3 partitions of the problems. Best per column in bold; underlined: within one SD of the best. Non-causal methods above the rule, causal below.}
\label{tab:language-detail-verdict}
\begin{tabular*}{\textwidth}{@{\extracolsep{\fill}}ll rrr rrr rrr@{}}
\toprule
 & & \multicolumn{3}{c}{Trained on Py} & \multicolumn{3}{c}{Trained on C++} & \multicolumn{3}{c}{Trained on Java} \\
\cmidrule(lr){3-5}\cmidrule(lr){6-8}\cmidrule(lr){9-11}
Model & Method & ID & $\to$C++ & $\to$Java & ID & $\to$Py & $\to$Java & ID & $\to$Py & $\to$C++ \\
\midrule
\multirow{4}{*}{Qwen3-Coder-30B} & Confidence & .608\,{\scriptsize$\pm$.028} & .637\,{\scriptsize$\pm$.052} & .687\,{\scriptsize$\pm$.011} & .637\,{\scriptsize$\pm$.054} & .611\,{\scriptsize$\pm$.026} & .688\,{\scriptsize$\pm$.013} & .689\,{\scriptsize$\pm$.010} & .605\,{\scriptsize$\pm$.027} & .630\,{\scriptsize$\pm$.053} \\
 & Probe & .670\,{\scriptsize$\pm$.011} & \underline{.628}\,{\scriptsize$\pm$.103} & .636\,{\scriptsize$\pm$.059} & \textbf{.874}\,{\scriptsize$\pm$.010} & .636\,{\scriptsize$\pm$.062} & .701\,{\scriptsize$\pm$.009} & \underline{.740}\,{\scriptsize$\pm$.044} & .670\,{\scriptsize$\pm$.026} & .739\,{\scriptsize$\pm$.020} \\
\cmidrule(lr){3-5}\cmidrule(lr){6-8}\cmidrule(lr){9-11}
 & Counterf.\ sim. & \underline{.671}\,{\scriptsize$\pm$.071} & \underline{.679}\,{\scriptsize$\pm$.129} & \underline{.721}\,{\scriptsize$\pm$.041} & \underline{.854}\,{\scriptsize$\pm$.021} & \underline{.714}\,{\scriptsize$\pm$.024} & .746\,{\scriptsize$\pm$.013} & \textbf{.751}\,{\scriptsize$\pm$.019} & \underline{.729}\,{\scriptsize$\pm$.017} & \textbf{.778}\,{\scriptsize$\pm$.013} \\
 & Value probe & \textbf{.717}\,{\scriptsize$\pm$.031} & \textbf{.692}\,{\scriptsize$\pm$.030} & \textbf{.737}\,{\scriptsize$\pm$.004} & .814\,{\scriptsize$\pm$.010} & \textbf{.731}\,{\scriptsize$\pm$.003} & \textbf{.761}\,{\scriptsize$\pm$.010} & \underline{.744}\,{\scriptsize$\pm$.011} & \textbf{.733}\,{\scriptsize$\pm$.019} & .712\,{\scriptsize$\pm$.026} \\
\midrule
\multirow{4}{*}{gpt-oss-20b} & Confidence & .542\,{\scriptsize$\pm$.023} & .529\,{\scriptsize$\pm$.016} & .549\,{\scriptsize$\pm$.050} & .528\,{\scriptsize$\pm$.016} & .541\,{\scriptsize$\pm$.023} & .549\,{\scriptsize$\pm$.050} & .555\,{\scriptsize$\pm$.045} & .525\,{\scriptsize$\pm$.009} & .525\,{\scriptsize$\pm$.016} \\
 & Probe & \underline{.812}\,{\scriptsize$\pm$.021} & \textbf{.692}\,{\scriptsize$\pm$.016} & .677\,{\scriptsize$\pm$.018} & \underline{.797}\,{\scriptsize$\pm$.011} & \underline{.651}\,{\scriptsize$\pm$.045} & .596\,{\scriptsize$\pm$.021} & .713\,{\scriptsize$\pm$.033} & \underline{.739}\,{\scriptsize$\pm$.055} & \underline{.651}\,{\scriptsize$\pm$.018} \\
\cmidrule(lr){3-5}\cmidrule(lr){6-8}\cmidrule(lr){9-11}
 & Counterf.\ sim. & \textbf{.828}\,{\scriptsize$\pm$.099} & \underline{.652}\,{\scriptsize$\pm$.144} & \textbf{.759}\,{\scriptsize$\pm$.034} & \textbf{.799}\,{\scriptsize$\pm$.088} & \textbf{.670}\,{\scriptsize$\pm$.168} & \textbf{.624}\,{\scriptsize$\pm$.057} & \textbf{.796}\,{\scriptsize$\pm$.027} & \textbf{.784}\,{\scriptsize$\pm$.092} & \textbf{.660}\,{\scriptsize$\pm$.091} \\
 & Value probe & .743\,{\scriptsize$\pm$.010} & .636\,{\scriptsize$\pm$.028} & .645\,{\scriptsize$\pm$.028} & .700\,{\scriptsize$\pm$.007} & .633\,{\scriptsize$\pm$.023} & .591\,{\scriptsize$\pm$.007} & .697\,{\scriptsize$\pm$.005} & .721\,{\scriptsize$\pm$.018} & \underline{.653}\,{\scriptsize$\pm$.015} \\
\bottomrule
\end{tabular*}
\end{table*}


% Numbers (2026-10-27): generated/tab_verdict_summary.tex,
% tab_language_detail_verdict.tex, paired_shift.csv (fold=mean, train=all: mean
% over the six language pairs, three partitions; "±" is 1.96 se_total). Per-pair
% statements are descriptive (fold 0/1/2 rows of paired_shift.csv, value probe
% − noncausal): ahead in all three partitions on pyctl→java, cpp→pyctl,
% cpp→java, java→pyctl; behind in all three on java→cpp. Within-problem AUROC
% from layer_selection{,_f1,_f2}.csv (mean over languages, then partitions);
% figure count from the score layer (artifact/manifest.csv): value-probe
% checkpoints above the probe's best layer on the shifted rows, Qwen3-Coder
% L35--L39 in each fold.

%The verdict target asks whether the model's judgment of a candidate can be
%trusted (Section~\ref{sec:problem}): the decision an agentic pipeline faces
%when it can either accept the verdict or run the tests. A high AUROC means that
%a method ranks candidates with a correct verdict above those with a wrong one.
Both LLMs are mediocre judges, with balanced accuracies of their verdicts of 57--74\% (Table~\ref{tab:datasets}) against the tests. RQ1 asks how well our methods identify the candidates on which the model's verdict is correct.
Table~\ref{tab:verdict-summary} summarizes the results, showing the mean per method over partitions.
Table~\ref{tab:language-detail-verdict} breaks down the language shift by training and evaluation language.

\paragraph{In distribution, hidden states outperform confidence}
On LCB and the language sets, the methods that read hidden states reach 0.69--0.83, while confidence is weak for
Qwen3-Coder (0.64--0.67) and close to chance for gpt-oss (0.54--0.57). %, as reported for natural-language tasks~\cite{Kadavath2022LanguageM, OrgadTGRSKB25}. 
On the language sets, the probe and the two causal methods are
level for Qwen3-Coder (0.758--0.761); for gpt-oss, counterfactual simulation is best (0.808). BCB is harder, no trained method exceeds 0.639 for Qwen3-Coder or 0.601 for gpt-oss.

\paragraph{Under language shift, the causal subspace is the most robust predictor for Qwen3-Coder}
%Probes on hidden states can degrade under distribution shift~\cite{OrgadTGRSKB25, Kramar2026Probes}. 
For Qwen3-Coder, the probe drops from 0.76 in distribution to 0.67 when the language changes%(mean over the six pairs)
, while both causal methods fall only from 0.76 to 0.73, against 0.64 for confidence. Averaged over the six language pairs, the value probe is ahead of
the best non-causal method of each pair by $0.042 \pm 0.021$ (95\% interval of the paired test, Section~\ref{sec:setup-protocol}) and ahead in each of the three partitions ($+0.035$, $+0.053$, $+0.038$). Counterfactual simulation
reaches the same mean but varies more between partitions ($\pm 0.043$). 
%Per pair, the value probe is ahead in all three partitions on four of the six pairs and behind in all three on Java$\to$C++, C++ being the hardest target language. 
The lead does not depend on the probe's layer choice: in 13 of the 18 pair--partition combinations, most value-probe checkpoints lie above the probe at its best layer on the shifted rows, as illustrated for Qwen3-Coder, trained on C++ in Figure~\ref{fig:verdict-profile}. 
%These results suggest that the learned subspace provides a more robust signal of verdict correctness than the full hidden state for Qwen3-Coder. 
For gpt-oss, all trained methods lose 0.07--0.12 AUROC under language shift, and neither causal method is significantly ahead of the best non-causal one. %Counterfactual simulation is ahead in five of the six pairs, but its lead varies between partitions from $-0.06$ to $+0.11$. 
The robustness reported for natural language tasks~\cite{Huang2025Internal} thus carries over to code for one of the two models; Section~\ref{sec:rq-subspace} examines what carries the gain.

% Figure: scripts/paper/fig_layer_profile.py --target verdict --cells qwen3:cpp
%   --wide --first-block 30 --out paper/generated/fig_verdict_profile_qwen3_cpp.pdf
% (fold 0, score layer; one row since 2026-10-28, probe below block 30 at
% 0.49--0.64 not drawn).
\begin{figure}[t]
\centering
\includegraphics[width=\columnwidth]{generated/fig_verdict_profile_qwen3_cpp}
\caption{AUROC per layer on the verdict target, Qwen3-Coder trained on C++
(fold 0), in distribution (left) and under language shift. Rings mark the
layer chosen in distribution; the dashed line is confidence. Under shift to
Python, every causal checkpoint lies above the probe's best layer; on Java,
most do.}
\label{fig:verdict-profile}
\end{figure}

% Figure: scripts/paper/make_result_tables.py --only fw (layer_selection{,_f1,_f2}.csv,
% columns pooled and within; 2026-10-27).
\begin{figure}[t]
\centering
\includegraphics[width=\columnwidth]{generated/fig_within}
\caption{Within-problem (hollow) against pooled AUROC (filled), in distribution
(faint) and under language shift: means over training languages and partitions.
Within-problem AUROC uses the problems with mixed outcomes only
(Section~\ref{sec:setup-protocol}); the length of a line is how much of the
pooled AUROC comes from ranking problems.}
\label{fig:within}
\end{figure}

\paragraph{The advantage lies in identifying unreliable problems, not candidates}
Pooled AUROC combines the ability to identify problems on which the judge is unreliable with the ability to distinguish correct from incorrect verdicts within a problem (Section~\ref{sec:setup-protocol}). The second is limited. Under language shift, within-problem AUROC is 0.616--0.618 for all trained methods of Qwen3-Coder. 
For gpt-oss the probe reaches 0.652, counterfactual simulation 0.635 and the value probe 0.576 (Figure~\ref{fig:within}).
The causal methods' pooled advantage thus comes mainly from identifying problems on which the judge is unreliable, rather than individual candidates it misjudges. 

\paragraph{A change of benchmark affects every method more than a change of
language}
Transfer between LCB %, whose problems are algorithmic, 
and BCB %, whose tasks rely on libraries, 
is asymmetric (Table~\ref{tab:verdict-summary}). From LCB to BCB, all trained methods of Qwen3-Coder fall
to 0.453--0.485, below confidence (0.552). In the reverse direction, the probe stays near chance (0.516), while the causal methods reach 0.680 and 0.695, close to confidence (0.669).
For gpt-oss, every method stays at 0.48--0.61 in both directions. %With one partition, these differences between methods are descriptive. 
%\orange{Section~\ref{sec:rq-subspace} relates the failure to the model's verdicts on BCB. }

\begin{rqanswer}{RQ1}
In distribution, hidden states predict whether the model's verdict is right far
better than its confidence. Under language shift, a probe on the causally
localized subspace is the most robust predictor for Qwen3-Coder (+0.04 AUROC over
the best non-causal method, in all three partitions), mainly by identifying problems
on which the judge is unreliable; for gpt-oss, no method is significantly ahead. 
Verdict prediction does not transfer reliably between benchmarks.
\end{rqanswer}


\subsection{RQ2: Which part of the causal methods carries their signal?}
\label{sec:rq-subspace}

On Qwen3-Coder's verdict target, the causal methods are the most robust predictors under language shift (RQ1). 
We now investigate which aspects of the method contribute to this performance: whether prediction requires the intervention itself, whether the learned directions matter beyond their dimensionality, and whether interchange intervention accuracy (see Section~\ref{sec:methods-causal}) indicates when the subspace is useful.


% Figure: scripts/paper/fig_subspace_layers.py (fold 0, results/*_randomlayers
% and the block-fixed learned value probes; 2026-10-22). Space permitting; it
% could replace Table~\ref{tab:subspace-control}.
\begin{figure}[t]
\centering
\includegraphics[width=\columnwidth]{generated/fig_subspace_layers}
\caption{Value probe on the learned DAS subspace against value probes on
three random 6-dimensional subspaces at the same layer, verdict target,
mean over shifted languages (the band is the range over
the random subspaces). Shaded: layers where DAS has
converged (IIA $\geq$ 0.9). On Qwen3-Coder the learned subspace is above
random at every layer; on gpt-oss it is above random before convergence and
near chance once DAS has converged.}
\label{fig:subspace-layers}
\end{figure}

\paragraph{Reading the subspace is sufficient}
Counterfactual simulation intervenes through the subspace and reads the answer
token, at the cost of $B$ additional forward passes per candidate; the value
probe reads the subspace directly in the judging pass. On Qwen3-Coder the two are equally robust under language shift (mean difference over the six pairs $-0.000 \pm 0.033$ on the verdict target, $+0.004 \pm 0.007$ on the program target), and both are ahead of confidence, which reads the same
answer token without the subspace ($+0.085 \pm 0.033$ and $+0.05 \pm 0.03$).
%Preliminary experiments showed that the choice of the $B=16$ bases barely matters: drawing them from different problems changed the pooled AUROC of one checkpoint by at most 0.01. 
The intervention thus brings no significant improvement over reading the
subspace, and the value probe costs no more at prediction time than a probe.

\paragraph{The learned direction matters, not the low dimension}
A six-dimensional readout might transfer better simply because it is small. We
therefore train the value probe on three random subspaces of the same
dimension at the same layer. At the layer chosen in distribution, the
learned subspace reaches 0.70--0.75 under shift for Qwen3-Coder, random
subspaces 0.48--0.61, no better than the probe on the full hidden state (0.668,
Table~\ref{tab:verdict-summary}). The same holds at every trained layer (Figure~\ref{fig:subspace-layers}). 
%\orange{Die Figure (insb. untere Zeile) wirft mehr Fragen auf... Auf Qwen3-Coder fokussierene evtl.?}

% % Undecided (2026-10-27): why the subspace helps on the verdict target;
% % hypotheses and checks in PAPER_PLAN_FORGE.md, open decision 14.
% \paragraph{\orange{The causal score reads how firmly a verdict is held}}
% \orange{The verdict score of counterfactual simulation is the probability that
% the bases output the candidate's verdict once its subspace is patched in. A
% base the patch does not move keeps its own answer, so the score carries an
% offset between \texttt{Yes} and \texttt{No} verdicts that depends on the bases,
% not on the candidate. Within one verdict the offset is constant, and there
% right verdicts score higher than wrong ones in every setting: under language
% shift, the AUROC among accepted and among rejected candidates is 0.65 and 0.77
% for Qwen3-Coder and 0.64 and 0.68 for gpt-oss (means over pairs and
% partitions). The causal methods thus read how firmly a candidate's state holds
% the model's decision. The offset explains the failure under benchmark shift:
% from LCB to BCB, Qwen3-Coder's \texttt{Yes} verdicts score 0.77 higher than its
% \texttt{No} verdicts, and 84\% of its wrong verdicts there are \texttt{Yes}, so
% the pooled AUROC falls to 0.43, although right verdicts still score higher within
% each verdict (0.57 and 0.61).}

\paragraph{Interchange accuracy does not guarantee predictive utility}
DAS learns the subspace so that swapping it between two inputs swaps the
model's verdict. IIA measures how often this succeeds. A high IIA thus shows
that the subspace controls the decision, not that it tells whether the decision
is right. On gpt-oss the two diverge (Figure~\ref{fig:subspace-layers}): at L11
and L13, where DAS succeeds only partly (IIA 0.54--0.64), the learned subspace
transfers better than random ones (0.56--0.68 against 0.52--0.55 under shift),
and L13 is the layer chosen in distribution. However, at L17--L21, where DAS almost
always succeeds (IIA 0.91--0.95), it falls to 0.52--0.58, no better than
random. %\orange{One reading is that the decision is fully formed at these layers, so the subspace holds a clean \texttt{Yes} or \texttt{No} that says little about whether it is right.} 
On Qwen3-Coder the converged layers (L39--L43, IIA 0.95) remain informative. IIA thus cannot tell the two cases apart, and layers are chosen by held-out prediction instead (Section~\ref{sec:setup-training}).

\begin{rqanswer}{RQ2}
Where the causal methods help, the learned subspace carries the gain: the value
probe, which reads it in the judging pass, is as robust as counterfactual
simulation, and random subspaces of the same size transfer only as well as the
full hidden state. %\orange{Within each verdict, right verdicts are held more firmly in the subspace than wrong ones; across verdicts, counterfactual simulation adds an offset set by its bases, which explains its failure from LCB to BCB.} 
IIA does not indicate when the subspace is useful: on gpt-oss, the
converged subspaces are no better than random ones.
\end{rqanswer}
 

\subsection{RQ3: How well do the methods predict whether the program is correct?}
\label{sec:rq-program}


The program target asks whether a candidate passes its tests, regardless of
the model's verdict. The two targets
agree only where the judge is right. %Across all candidates, they are unrelated: the model judges passing and failing programs equally accurately. %(Adrian: die zwei Sätze mag ich nicht. Sicher dass korrekt?)}
We therefore ask whether the results on the verdict target carry over. Table~\ref{tab:program-summary} summarizes the results. In distribution, all methods predict
program correctness well above chance, AUROC ranges from 0.70--0.88 on LCB and the language
sets, with the hidden-state probe performing best on the language sets for both models. BCB is more challenging for every method.

% Generated by scripts/paper/make_result_tables.py -- do not edit.
\begin{table}[t]
\centering\footnotesize
\setlength{\tabcolsep}{3.2pt}
\caption{AUROC for program correctness (mean over 3 partitions). \emph{LCB}, \emph{BCB}: trained and evaluated on that benchmark; $\to$: evaluated on the other one. Language shift: \emph{ID} is the mean over the three training languages, \emph{Shift} the mean over the six train$\to$evaluation pairs, $\Delta$ their difference. %In-distribution cells are cross-fitted over problems; shifted cells use the layer chosen in distribution. 
Best per column and model in bold.}
\label{tab:program-summary}
\begin{tabular}{l rrrr rrr}
\toprule
 & \multicolumn{4}{c}{Benchmark shift} & \multicolumn{3}{c}{Language shift} \\
\cmidrule(lr){2-5}\cmidrule(lr){6-8}
Method & LCB & $\to$BCB & BCB & $\to$LCB & ID & Shift & $\Delta$ \\
\midrule
\multicolumn{8}{l}{\emph{Qwen3-Coder-30B}} \\
Confidence & .815 & .650 & .648 & .815 & .759 & .755 & \textbf{$-$.003} \\
Probe & \textbf{.844} & \textbf{.656} & \textbf{.688} & .697 & \textbf{.878} & .790 & $-$.088 \\
LAT & .844 & .580 & .590 & .731 & .807 & .799 & $-$.009 \\
\addlinespace[2pt]
Counterf.\ sim. & .839 & .644 & .635 & .813 & .848 & .806 & $-$.042 \\
Value probe & .838 & .642 & .664 & \textbf{.815} & .861 & \textbf{.810} & $-$.051 \\
\midrule
\multicolumn{8}{l}{\emph{gpt-oss-20b}} \\
Confidence & .723 & \textbf{.648} & .646 & .721 & .696 & .691 & \textbf{$-$.004} \\
Probe & .793 & .575 & .649 & .690 & \textbf{.859} & .693 & $-$.166 \\
LAT & .767 & .618 & .608 & .676 & .760 & .666 & $-$.094 \\
\addlinespace[2pt]
Counterf.\ sim. & .780 & .647 & .652 & \textbf{.724} & .820 & \textbf{.710} & $-$.111 \\
Value probe & \textbf{.803} & .638 & \textbf{.652} & .712 & .846 & .666 & $-$.180 \\
\bottomrule
\end{tabular}
\end{table}



\paragraph{Under language shift, the causal methods lose their advantage}
The full-state probe loses 0.088 AUROC for Qwen3-Coder and 0.166 for gpt-oss, while confidence changes little ($-0.003$ and $-0.004$). For Qwen3-Coder, LAT and both causal methods nevertheless remain significantly ahead of confidence by 0.044--0.054 AUROC on average; the probe's lead of 0.035 is no longer significant. For gpt-oss, no trained method significantly outperforms confidence. Importantly, neither causal method outperforms the best non-causal method: for Qwen3-Coder, the value probe and counterfactual simulation trail it by 0.010 and 0.014 AUROC, respectively, neither significantly. For gpt-oss, the value probe is significantly behind ($-0.068$), while counterfactual simulation is statistically level ($-0.025$). Thus, the robustness of causal localization observed for verdict correctness (RQ1) does not carry over to program correctness.


% In distribution, the trained methods are 0.05--0.16 ahead of confidence on the
% language sets. Under language shift the probe loses most ($-0.088$ for
% Qwen3-Coder, $-0.166$ for gpt-oss), while confidence barely changes ($-0.003$,
% $-0.004$). For Qwen3-Coder, LAT and the two causal methods remain
% significantly ahead of confidence, by 0.044--0.054 on average over the pairs;
% the probe's lead ($+0.035 \pm 0.035$) is no longer significant. For gpt-oss,
% no trained method is ahead ($-0.037$ to $+0.019$, none significant). Where no labeled data exists for the new language, confidence is thus about 0.05 behind the best
% trained method for Qwen3-Coder and as good as any for gpt-oss.

% \paragraph{The causal methods have no advantage on program correctness}
% On average over the pairs, neither causal method is ahead of the best
% non-causal method under language shift: for Qwen3-Coder, the differences are
% $-0.010 \pm 0.019$ (value probe) and $-0.014 \pm 0.017$ (counterfactual
% simulation), neither significant; for gpt-oss, the value probe is
% significantly behind ($-0.076 \pm 0.035$) and counterfactual simulation is
% level ($-0.025 \pm 0.034$). Their robustness on the verdict target (RQ1) does
% not carry over to program correctness.

\paragraph{Program correctness transfers between benchmarks}
In contrast to verdict correctness, predictors for program correctness transfer between LCB and BCB in both directions. When trained on LCB and evaluated on BCB, methods reach 0.580--0.656 for Qwen3-Coder, and 0.575--0.648 for gpt-oss, comparable to
the respective in-distribution results on BCB (0.590--0.688 and 0.608--0.652).
 Transfer in the reverse direction is similarly strong, reaching 0.70--0.82 for Qwen3-Coder and 0.68--0.72 for gpt-oss.
% Figure removed for space (2026-10-23): the selection cost is given in the
% text. scripts/paper/fig_layer_profile.py --target program --cells qwen3:cpp
% gptoss:cpp
% \begin{figure}[t]
% \centering
% \includegraphics[width=\columnwidth]{generated/fig_program_profile_qwen3_cpp_gptoss_cpp}
% \caption{AUROC per layer on the program target, in distribution (top) and under
% language shift (fold 0), trained on C++: Qwen3-Coder (left), gpt-oss
% (right). Rings mark the layer chosen in distribution; the dashed line is
% confidence at its chosen temperature. The layer that is best for the probe in
% distribution is not the one that transfers. LAT is omitted: its per-layer sign
% fit flips in many layers.}
% \label{fig:layers}
% \end{figure}
\paragraph{Pooled AUROC overstates candidate-level discrimination}
As for verdict correctness, pooled AUROC reflects both the ability to identify problems the model tends to solve and the ability to distinguish candidates within a problem.
Choosing a candidate needs only the second, and here every method is weak (Figure~\ref{fig:within}, bottom): under language shift, within-problem AUROC lies at 0.58--0.64 for all methods and both models, against a pooled AUROC of 0.67--0.81. For Qwen3-Coder confidence performs best within problem (0.629). 
The trained methods reach higher within-problem AUROCs of 0.69--0.78 only for gpt-oss in distribution, but lose this advantage under shift. Consequently, the predictors are more useful for identifying problems that warrant verification than for choosing among candidates of the same problem.
Section~\ref{sec:rq-problem} predicts at problem level directly.

% Numbers below: generated/tab_genprobe.tex from generated/genprobe_compare.csv
% (scripts/analysis/genprobe_compare.py, fold 0, 2026-10-22).
\paragraph{Reading the generation is as good as asking}
All methods so far read the model in a separate judging pass. The same information might already be present when the model generates the candidate. %\orange{Verstehe ich nicht. Ist das nicht per Definition identisch?}. 
We therefore train the same probe on the hidden states of the model's own
generation (without the judgment prompt), replayed token by token, and compare it with the judging-pass probe on identical candidates, Table~\ref{tab:genprobe} reports the results. We compare different token positions here: at the last token of the candidate's code, it is
at least as good as the judging pass in distribution (0.895 against 0.886 for
Qwen3-Coder, 0.864 against 0.851 for gpt-oss); under language shift, it is
close for Qwen3-Coder (0.795 against 0.779) and behind for gpt-oss (0.673
against 0.704). A correctness score can thus be read while the code is
written. Read at the last prompt token, before any code exists, the probe
has seen only the problem, yet for Qwen3-Coder it matches the judging pass in
distribution (0.891). Much of what the methods predict is thus the problem,
not the candidate.

% Generated by scripts/paper/make_result_tables.py -- do not edit.
\begin{table}[t]
\centering\footnotesize
\setlength{\tabcolsep}{3.2pt}
\caption{The probe on the judging pass against the same probe on the states of the model's own generation, program target on the language sets. \emph{ID}: mean AUROC over the three languages, \emph{Shift}: mean over the six language pairs, \emph{Within}: AUROC over pairs of candidates of the same problem. The last prompt token has seen only the problem.}
\label{tab:genprobe}
\begin{tabular}{l rr rr}
\toprule
 & \multicolumn{2}{c}{ID} & \multicolumn{2}{c}{Lang. Shift} \\
\cmidrule(lr){2-3}\cmidrule(lr){4-5}
Readout & Pooled & Within & Pooled & Within \\
\midrule
\multicolumn{5}{l}{\emph{Qwen3-Coder-30B}} \\
Judging pass & .886 & .553 & .779 & .540 \\
Generation, last code token & .895 & .620 & .795 & .627 \\
Generation, last prompt token & .891 & .514 & .755 & .531 \\
\midrule
\multicolumn{5}{l}{\emph{gpt-oss-20b}} \\
Judging pass & .851 & .772 & .704 & .609 \\
Generation, last code token & .864 & .781 & .673 & .623 \\
Generation, last prompt token & .720 & .524 & .669 & .478 \\
\bottomrule
\end{tabular}
\end{table}


\begin{rqanswer}{RQ3}
All methods predict program correctness well above chance, largely by recognizing the problems the model fails on, %Within one problem, every method is weak 
and this transfers better between benchmarks than the verdict. Under
language shift, the trained methods lose most of their lead over confidence:
about 0.05 remains for Qwen3-Coder, nothing for gpt-oss, and the causal
methods' robustness on the verdict target does not carry over. A probe on the model's own generation is as good as one on a separate judging pass in distribution.
\end{rqanswer}


\subsection{RQ4: Does agreement across samples improve problem-level prediction?} % tell better than a single pass which problems the model will solve?}
\label{sec:rq-problem}
% Numbers from generated/tab_problem_level.tex and generated/problem_level_ci.csv
% (scripts/analysis/problem_level_ci.py, 2026-10-16). "Significant": 95% paired
% problem bootstrap (2000 draws, same rows and averaging as the table) excludes
% zero. Dropped 2026-10-28: the before-sampling paragraph (problem-only probe
% now in RQ3), the end-of-reasoning probe, per-pair counts, the Harrell's C
% note, the sentence on prior consistency studies (covered in related work).

This question moves from the candidate to the problem level. Much of the
candidate-level signal reflects which problems the model solves
(Section~\ref{sec:rq-program}). When several candidates are sampled for a
problem, their agreement provides a further signal.
EigenScore~\cite{Chen24Eigenscore} reads this agreement from the hidden states
of the generations (Section~\ref{sec:methods-consistency}). Like confidence, it
needs no training, so it cannot lose under shift. We compare at equal sampling
cost: each single-pass method is averaged over the candidates EigenScore uses,
and every method predicts whether at least half of a problem's candidates pass
(Table~\ref{tab:problem-level}).

% - A possibility to prevent the problems under distribution shift, is to use an unsupervised method, that needs no training. Also other advantages, no labeled training data needed\dots
% One such method is eigenscore, which uses cnsistency across the hidden states of n different samples (n=5 in our experiments) to compute a score for 'how likely that this problem will be solved by the model'. As a consistency-based method, eigenscore scores at problem level in contrast to our other candidate level methods. To compare it to our methods nonetheless, we mean-pool our candidate scores to problem scores and label them by majority vote -> a problem is solved if more than $50\%$ of its candidates pass (reasonable since most (X-Y\% depending on the benchmark) of our problems contain passing or failing candidates only).

% Remember that we exclude empty generations (where no program could be extracted from the models reasoning) as trivial from our experiments, so we filter out all problems with less than n=5 non-empty samples, resulting in a slightly reduced test set for this experiment.

% Generated by scripts/paper/make_result_tables.py -- do not edit.
\begin{table}[t]
\centering\footnotesize
\setlength{\tabcolsep}{3.2pt}
\caption{AUROC for whether a problem is solved (at least half of its candidates pass). Single-candidate methods are averaged over the candidates of a problem, at the layer of Table~\ref{tab:program-summary}; EigenScore is computed from the sampled candidates and needs no training. Columns as in Table~\ref{tab:program-summary}. $\uparrow$: significantly above EigenScore (paired problem bootstrap, 95\%; for ID and Shift on the mean over languages). All methods on the problems that have an EigenScore.$^{*}$}
\label{tab:problem-level}
\begin{tabular}{l rrrr rr}
\toprule
 & \multicolumn{4}{c}{Benchmark shift} & \multicolumn{2}{c}{Language shift} \\
\cmidrule(lr){2-5}\cmidrule(lr){6-7}
Method & LCB & $\to$BCB & BCB & $\to$LCB & ID & Shift \\
\midrule
\multicolumn{7}{l}{\emph{Qwen3-Coder-30B}} \\
EigenScore & .504 & .658 & .658 & .504 & .528 & .528 \\
Confidence & .877$^{\uparrow}$ & \textbf{.670} & .671 & .872$^{\uparrow}$ & .774$^{\uparrow}$ & .775$^{\uparrow}$ \\
Probe & .906$^{\uparrow}$ & .638 & .671 & .884$^{\uparrow}$ & \textbf{.911}$^{\uparrow}$ & .803$^{\uparrow}$ \\
LAT & .902$^{\uparrow}$ & .618 & \textbf{.682} & .871$^{\uparrow}$ & .844$^{\uparrow}$ & .821$^{\uparrow}$ \\
\addlinespace[2pt]
Counterf.\ sim. & \textbf{.912}$^{\uparrow}$ & .646 & .676 & .881$^{\uparrow}$ & .875$^{\uparrow}$ & .831$^{\uparrow}$ \\
Value probe & .904$^{\uparrow}$ & .652 & .680 & \textbf{.889}$^{\uparrow}$ & .890$^{\uparrow}$ & \textbf{.838}$^{\uparrow}$ \\
\midrule
\multicolumn{7}{l}{\emph{gpt-oss-20b}} \\
EigenScore & .660 & .518 & .518 & .660 & .586 & .586 \\
Confidence & .698 & \textbf{.677}$^{\uparrow}$ & .678$^{\uparrow}$ & .692 & .623 & .615 \\
Probe & .791 & .615 & \textbf{.712}$^{\uparrow}$ & \textbf{.759} & \textbf{.845}$^{\uparrow}$ & \textbf{.676}$^{\uparrow}$ \\
LAT & \textbf{.808}$^{\uparrow}$ & .662$^{\uparrow}$ & .622$^{\uparrow}$ & .688 & .772$^{\uparrow}$ & .645 \\
\addlinespace[2pt]
Counterf.\ sim. & .772 & .655$^{\uparrow}$ & .685$^{\uparrow}$ & .692 & .800$^{\uparrow}$ & .662 \\
Value probe & .790 & .663$^{\uparrow}$ & .692$^{\uparrow}$ & .676 & .834$^{\uparrow}$ & .630 \\
\bottomrule
\end{tabular}
\par\smallskip
\parbox{\linewidth}{\scriptsize $^{*}$EigenScore needs $m$ non-empty samples: five per problem on LCB and the language sets, ten on BCB. Problems with fewer non-empty generations are left out for every method, and the single-candidate methods are averaged over the candidates of the same samples. This keeps 93--100\% of the evaluation problems for Qwen3-Coder-30B and 76--94\% for gpt-oss-20b. %Of the problems left out (counted once per benchmark), 0 of 40 are solved for Qwen3-Coder-30B and 47 of 143 for gpt-oss-20b.
}
\end{table}


% Numbers: generated/tab_genprobe.tex (candidate-level pooled AUROC, fold 0).
% A problem-level row for the last-prompt-token probe in
% Table~\ref{tab:problem-level} would make this directly comparable with
% EigenScore (make_result_tables.py t5); not done (2026-10-28).
The cheapest prediction needs no sample at all: for Qwen3-Coder, a probe that
has seen only the problem %\orange{verstehe ich nicht} 
already ranks candidates as well as the judging pass
(Section~\ref{sec:rq-program}). Sampling several candidates pays off for
prediction only if their agreement adds to what a single pass provides.

\paragraph{Agreement provides no consistent advantage}
For Qwen3-Coder, EigenScore is at chance on LCB (0.504) and on the language
sets (0.528), where every single-candidate method, confidence included, is
significantly ahead (0.77--0.91). On BCB it is comparable to them (0.658 against
0.671--0.682). For gpt-oss the benchmarks swap roles: EigenScore is at chance
on BCB (0.518 against 0.622--0.712) and closer on LCB (0.660 against
0.698--0.808). Under shift, the single-pass methods lose ground. 
From LCB to BCB, Qwen3-Coder's trained methods fall slightly below EigenScore (0.618--0.652 against 0.658) while confidence stays above it (0.670); under language shift, the
gap remains large for Qwen3-Coder (0.775--0.838 against 0.528) and narrows for
gpt-oss (0.615--0.676 against 0.586).


\begin{rqanswer}{RQ4}
Agreement across samples gives no consistent advantage over a single pass.
EigenScore is level with the single-pass methods in some cells and near
chance in others, and confidence, which needs no labels either, is never
significantly below it. That EigenScore cannot lose under shift does not help where it starts near chance.
\end{rqanswer}

% Notes from earlier drafts:
% RQ2: How well can they tell whether the model's own verdict on its code can
% be trusted, and how much of that survives shift?
% How does an unsupervised consistency-based method compare to these methods?
% Numbers are reported in mean±std. For each task, the highest AUROC value is
% in bold. Values that are within a standard deviation are underlined.
